\documentclass[11pt]{article}
\usepackage[margin=1.1in]{geometry}
\usepackage{amsmath,amssymb,amsthm,booktabs}
\usepackage[hidelinks]{hyperref}

\newtheorem{theorem}{Theorem}
\newtheorem{proposition}[theorem]{Proposition}
\theoremstyle{remark}
\newtheorem{remark}[theorem]{Remark}

\newcommand{\vp}{\varpi}
\newcommand{\eps}{\varepsilon}
\newcommand{\Sh}{\circledcirc}
\newcommand{\lean}[1]{\texttt{#1}}
\newcommand{\dA}{d_{A_0}}

\title{Bug Fixes for Geometric Unity:\\ A Charitable Formalization in Lean}
\author{Claude\thanks{Claude is an AI model made by Anthropic. The B~Man chose the topic and
directed the work; Claude wrote the Lean code and this note. It is not a statement by
Anthropic. Every mathematical claim marked as checked is verified by the Lean kernel and does
not depend on trusting either author.} \and The B~Man}
\date{November 2026}

\emergencystretch=3em
\begin{document}
\maketitle

\begin{abstract}
We formalize in Lean~4 with Mathlib as much of E.~Weinstein's \emph{Geometric Unity: Author's
Working Draft v1.0} as is stated precisely enough to formalize, reading the draft as charitably
as we can. The finite-dimensional claims of Sections~3, 4 and~8 are correct. A decomposition
that the draft asks to have checked is consistent. Sections~6 and~7 contain a sign error that
makes two lemmas false as printed, and no single sign convention makes the printed lemmas
simultaneously true; we give a three-sign repair under which all of them hold. After the
repair, and granting every standard fact the proof assistant's library cannot yet express, one
thing remains: the field equation is built from an operator the draft does not define. We then
try to supply the missing definition from the literature. With the textbook choice of forms,
the draft's sample operator is the Ricci tensor if read as printed, and is exactly the Einstein
tensor, with the draft's coefficient, under one of eight readings of its brackets. Finally we
count the operators that the draft's description allows. There are exactly $200$ independent
ones, every one of them annihilates the Weyl curvature, and gravity constrains four of their
parameters. For none of them except zero can the field equation be the Euler--Lagrange equation
of an action in the way the draft states. So the part of the operator that can be reconstructed
is the part that was already known, and the part that would be new is missing and cannot be
supplied in the form the draft describes.
\end{abstract}

\section{Ground rules}

The draft \cite{GU} says of itself that ``the interested community is welcome to send bug
fixes'' (footnote~9). This note takes that literally. A proof assistant is a good tool for the
job because it has no opinion of the author. It accepts what is defined and proved and stops at
what is not. We use Lean~4 \cite{Lean4} and its mathematical library Mathlib \cite{Mathlib}.
Every quotation from the draft below was compared with the PDF cited as \cite{GU}, and
equation, lemma and footnote numbers are the draft's own.

The draft also says, in the footnote to its title, that its author ``is not a physicist and is
no longer an active academician, but is an Entertainer'', and it calls itself ``this work of
entertainment''. We do not take that literally. The document states definitions, lemmas and
proofs, and it is offered in public as a theory of physics, so we check it as mathematics. A
lemma is true or false whatever its author calls the document it appears in. The same footnote
asks that the work not be built upon without permission; this note builds nothing on it and
only checks what is there.

We read the draft under four rules, each chosen in its favor.

\begin{enumerate}
\item \textbf{Standard mathematics is granted.} The revision of Mathlib we build against has
Clifford algebras and the spin group, a Hodge star on the exterior powers of a single vector
space, and covariant derivatives on vector bundles, including the Levi-Civita connection. It has
no principal bundles, no spin structures or spinor bundles, no connections on associated
bundles or their curvature, and no Hodge star on bundle-valued forms. Wherever the draft uses
such things in the standard way \cite{KN,LM}, we assume they exist and behave as the textbooks
say. None of
Mathlib's gaps is counted against the draft.
\item \textbf{Conflicting conventions are resolved in the draft's favor.} Where the text is
inconsistent we look for the reading that makes the most statements true.
\item \textbf{What can be repaired is repaired.} An error with a fix is reported with the fix.
\item \textbf{No physics is judged.} We ask only whether the mathematics is defined and whether
the stated lemmas follow. The objections of Nguyen and Polya concerning the chiral anomaly,
unitarity and supersymmetry \cite[Sections~3.2 and~3.3]{NP} are outside what a proof assistant
checks, and we take no position on them.
\end{enumerate}

Sections~5--7 of the draft are algebra on three pieces of data, and the draft's own proofs use
nothing else: a group $\mathcal H$ (the gauge group), an abelian group $\mathcal N$ on which it
acts (the ad-valued one-forms, with $h\cdot n = h n h^{-1}$), and a map
$c:\mathcal H\to\mathcal N$, $c(h)=h^{-1}(\dA h)$, satisfying the Leibniz rule
$c(h_1h_2)=h_2^{-1}c(h_1)h_2+c(h_2)$. In the geometric setting these are the group of gauge
transformations, the one-forms with values in the adjoint bundle, and the change of the
connection $A_0$ under a gauge transformation \cite[Section~2.1]{DK}. A map with this Leibniz
property is a crossed homomorphism, or $1$-cocycle, for the right action of $\mathcal H$ on
$\mathcal N$, and it is a classical fact that crossed homomorphisms correspond to the
homomorphisms $h\mapsto(h,c(h))$ into the semidirect product that split the projection onto
$\mathcal H$ \cite[Chapter~IV, Section~2]{Brown}. The draft's Lemma~6.1 is an instance of that
fact. We formalize at that level of generality. A statement
proved from these axioms holds in the draft's geometric setting. A statement that fails from
these axioms fails there too whenever the explicit error term is nonzero.

\section{What checks}

Table~\ref{tab} lists every statement examined. Everything before Section~6 that could be
checked is correct.

\begin{table}[h]
\centering\small
\begin{tabular}{@{}lll@{}}
\toprule
Draft & Statement & Result \\
\midrule
\S3.1, (3.8)--(3.9) & Signature of the Frobenius form on the fibre of metrics & True \\
\S3.5 & Signature totals leading to $(7,7)$ & True \\
(3.21)--(3.27) & Clifford degrees forming $gl(64,\mathbb R)$, $so(64,64)$, $u(64,64)$ & Consistent \\
(4.2)--(4.3) & Internal spinor dimensions $32$ and $16$ & True \\
(8.2)--(8.6) & $2^{14}=128^2$ and the degree split of $u(64,64)$ & True \\
Def.~5.1, (5.9) & The inhomogeneous gauge group $\mathcal G$ is a group & True \\
(6.2) & Right action of $\mathcal G$ on connections & True \\
Lemma 6.1 & The tilted map $\tau$ is a homomorphism & True \\
\S6.2 & $\hat\tau$ parameterizes the stabilizer of $A_0$ & True \\
Lemma 6.3 & $\pi$ is constant on tilted orbits & True as printed \\
Lemma 6.4 & Left action of $\mathcal G$ on $\mathcal N$ & True \\
Lemma 6.5 & The bi-connection map is equivariant & \textbf{False as printed} \\
Lemma 7.2 & The augmented torsion is equivariant & \textbf{False as printed} \\
\S\S6--7 & Lemmas 6.3 and 7.2 under one sign of $\tau$ & \textbf{Inconsistent} \\
Def.~8.1 & The forms $\Phi_i$ & \textbf{Not determined} \\
(9.2) & The shiab operator & \textbf{Not defined} \\
(8.1) & The conjugation recipe is equivariant for every $\Phi$ & True \\
(9.3) & The sample operator is Einstein's contraction & One reading of eight \\
(9.8) & Gauge covariance of the field equation & Conditional \\
Def.~8.1 & Dimension of each space the $\Phi_i$ are drawn from & Exactly $2$ \\
(9.2) & Number of independent operators of the stated kind & $200$ \\
\S9.1, \S12.4 & The operator annihilates the Weyl curvature & True for all $200$ \\
(9.4), (9.7) & (9.8) is the Euler--Lagrange equation of the action & \textbf{False for every operator} \\
\bottomrule
\end{tabular}
\caption{Statements examined. All results are machine-checked except the row for (9.3), which is numerical, the last four rows, whose status is given in Sections~\ref{sec:count} and~\ref{sec:action}, and where Section~\ref{sec:limits} says otherwise.}
\label{tab}
\end{table}

Two of these deserve a remark.

\paragraph{The signature.} For a metric of signature $(1,3)$ the form (3.8),
$\langle A,B\rangle=\operatorname{Tr}(g^{-1}Ag^{-1}B)$, is diagonal on symmetric matrices with
seven coefficients of one sign and three of the other, and the trace direction lies among the
seven. So (3.8) as defined has signature $(7,3)$, or $(3,7)$ with the overall sign for which
the draft's traceless part is $(3,6)$, and the $(4,6)$ of (3.9) comes from then reversing the
sign of the trace direction by hand. The draft says this is a choice, in the text before (3.9)
and in footnote~2, so we record it as true.

\paragraph{A request fulfilled.} Footnote~3 says of the decomposition (3.27) of $u(64,64)$
into exterior powers that it ``was taken from a decades old file and should be checked by
someone more current''. We checked it. The rule for the degrees appearing in a Clifford
commutator is standard: the product of elements of degrees $p$ and $q$ has components in
degrees $p+q-2k$, and exchanging the factors multiplies the component of degree $p+q-2k$ by
$(-1)^{pq-k}$. This follows from the fact that reversal of the order of products is an
anti-automorphism of the Clifford algebra acting on degree $m$ by $(-1)^{m(m-1)/2}$
\cite[Chapter~I, Section~1]{LM}. Given this rule, the degrees congruent to $1,2 \bmod 4$ are closed under the bracket, the
degrees congruent to $0,3\bmod 4$ twisted by $i$ bracket with them and with each other as
(3.27) requires, and the dimensions are $8128+8256=128^2$. The decomposition is consistent.

\section{The bug}

The draft defines the tilted map in (6.4) by
\[
\tau(h) = \bigl(h,\; -h^{-1}(\dA h)\bigr) = (h, -c(h)),
\]
and the augmented torsion in (7.3) by $T_{(\eps,\vp)}=\vp-c(\eps)$. Lemma~7.2 asserts
$T_{g\cdot\tau(h)}=h^{-1}T_g h$. Using the group law (5.9) and the Leibniz rule,
\[
T_{g\cdot\tau(h)}
= \bigl(h^{-1}\vp h - c(h)\bigr) - \bigl(h^{-1}c(\eps)h + c(h)\bigr)
= h^{-1}T_g h \;-\; 2\,h^{-1}(\dA h).
\]
The draft's proof (7.4) writes $-h^{-1}(\dA h)$ in its first line and $+h^{-1}(\dA h)$ in its
second. The same change of sign occurs between the lines of (6.20), the proof of Lemma~6.5.

\begin{theorem}[\lean{torsion\_mul\_tauMinus}, \lean{biConn\_mul\_tauMinus}]
With $\tau$ as printed in \textup{(6.4)}, Lemma~7.2 fails by $-2c(h)$, and the first component
of Lemma~6.5 fails by the same term.
\end{theorem}

The charitable response is that the sign in (6.4) is a misprint. With
$\tau(h)=(h,+c(h))$, which is also the form transcribed in \cite[Section~2.2]{NP} from the
video presentation \cite{Video},
Lemmas~6.5 and~7.2 are true. But Lemma~6.3, with $\pi$ as printed in (6.11), then fails:
\[
\pi\bigl(g\cdot\tau(h)\bigr) = \pi(g) + 2\,\eps\,(\dA h)\,h^{-1}\eps^{-1}.
\]

\begin{theorem}[\lean{lemma63\_and\_lemma72\_incompatible}]
Suppose $\mathcal N$ has no $2$-torsion. Fix either sign for $\tau$. If Lemma~6.3 and Lemma~7.2
both hold as printed for that sign, then $h^{-1}(\dA h)=0$ for every $h\in\mathcal H$.
\end{theorem}

So the printed text cannot be made true by choosing a convention. The draft anticipates
something of the kind: footnote~7 says there ``may be 2 or more conflicting convetions [\emph{sic}] in this section''. The theorem says where they are and that they cannot be reconciled without changing
a definition.

\section{The fix}

Change three signs and leave everything else as printed:
\begin{align*}
\text{(6.4)}\quad & \tau(h) = \bigl(h,\; +h^{-1}(\dA h)\bigr),\\
\text{(6.11)}\quad & \pi(\eps,\vp) = \eps\vp\eps^{-1} - (\dA\eps)\eps^{-1},\\
\text{(6.15)}\quad & (\eps,\vp)\cdot\gamma = \eps(\vp+\gamma)\eps^{-1} - (\dA\eps)\eps^{-1}.
\end{align*}

\begin{theorem}[\lean{GU.Repaired}]
With these three changes, Lemmas~6.1, 6.3, 6.4, 6.5 and~7.2 all hold, the stabilizer of $A_0$
is the image of the map $\hat\tau$ of \textup{(6.9)} as printed, the action \textup{(6.15)} is
the one induced by $\pi$ from left multiplication, and right multiplication by $\tau(h)$ acts on
$A_0+\vp$ as the gauge transformation $h$.
\end{theorem}

The repair also resolves an oddity in the printed text, where the ``tilted'' map $\tau$ of
(6.4) and the ``flipped'' map $\hat\tau$ of (6.9) are given by the same formula. After the
repair they are different maps with different jobs, which is evidently what was meant.

We regard Sections~5--7, so repaired, as correct. From here on we use the repaired version.

\section{What cannot be fixed}

The first-order field equation (9.8) is
\[
\Sh_\eps\, F_{A_\omega} \;=\; -\,\kappa_1 *T_\omega .
\]
The curvature $F$ and the Hodge star are standard \cite{KN,LM} and we grant them. $T$ is the torsion of
Section~7, now verified. The remaining ingredient is the shiab operator
$\Sh_\eps:\Omega^2(Y,\mathrm{ad})\to\Omega^{d-1}(Y,\mathrm{ad})$. It is the one new operator
in the equation, and the draft says the following about it.

\begin{itemize}
\item ``Then a Shiab contraction operator might look like: [(8.1)] or perhaps a sum of such
terms.''
\item ``The author is no longer in a position to go chasing after the complete picture and
simply details some of the available tools for customizing such operators.''
\item ``The particulars of the Shiab operators are workmanlike and not of much interest.''
\item ``Unfortunately, the author is no longer conversant in that language and has been unable
to locate the notes from decades ago that originally picked out the operator of choice.''
\item ``The Shiab operator that he settled on (but cannot yet now locate) was chosen for its
properties relative to the Bianchi identity.'' (footnote~10)
\end{itemize}

These are from Sections~8, 8.2 and~9.1 of the draft. A repair requires something to repair.
Here the charitable reading has nothing to work with, because the object is absent rather than
wrong.

We are not the first to stop at this operator. Nguyen and Polya objected to it in writing
before the draft existed \cite[Section~3.1]{NP}. Their objection is to the version in the video
\cite{Video}, where the Lie algebra is $u(128)$ and the definition needs a complexification
that is not stated. The draft appeared five weeks after their response, with $u(64,64)$ and the
complexified Clifford algebra in (8.4)--(8.6), which are among the statements checked in
Section~2. What we add is that in the written draft the operator is still not specified, and,
in Sections~\ref{sec:count} and~\ref{sec:action}, a count of what the specification would have
to choose among.

The draft does offer one sample operator, (9.3), ``which makes the parallel to Einstein's
contraction of the full Riemannian curvature explicit''. We tried to take that as the
definition. It is built from forms $\Phi_1$ and $\Phi_2$, and Definition~8.1 says only that the
$\Phi_i$ are ``a basis for the invariant subspaces of''
$[\Lambda^i(\mathbb R^{7,7})\otimes u(64,64)]^{Spin(7,7)}$. This does not determine them.

\begin{proposition}[\lean{equivariant\_mul\_commuting}]
Let a group act on a ring $A$ by conjugation and on a set $V$. If $\Phi:V\to A$ is equivariant
and $\lambda\in A$ commutes with the group, then $v\mapsto\Phi(v)\lambda$ is equivariant.
\end{proposition}

In $Cl(7,7)$ the volume element $\lambda$ anticommutes with vectors, because the dimension is
even, and so commutes with the even subalgebra and with $Spin(7,7)$ inside it
\cite[Chapter~I, Section~3]{LM}. Taking $\Phi(v)=\gamma(v)$
gives a second invariant $\gamma(v)\lambda$, of Clifford degree $13$ where the first has
degree~$1$. The draft's own equation (8.5) places both degrees inside $so(64,64)\subset
u(64,64)$. The same holds for $\Phi_2$ with degrees $2$ and~$12$, by (8.5) and (8.6). So each
space the $\Phi_i$ are drawn from is at least two-dimensional (Section~\ref{sec:count} shows
that it is exactly two-dimensional), and the sample operator (9.3)
has free parameters that the draft does not fix, in addition to the constant $\kappa_1$, which
is given no value.

What can be stated in Lean is therefore this.

\begin{theorem}[\lean{upsilon\_equivariant}]\label{thm:upsilon}
Let $\Sh$ be \emph{any} family of additive maps
$\Sh_\eps:\Omega^2\to\Omega^{d-1}$ satisfying
$\Sh_{\eps h}(h^{-1}\xi h)=h^{-1}(\Sh_\eps\xi)h$. Then, granting the standard equivariance of
curvature and of the Hodge star,
$\Upsilon_g=\Sh_\eps F_g+\kappa_1 *T_g$ satisfies
$\Upsilon_{g\cdot\tau(h)}=h^{-1}\Upsilon_g h$.
\end{theorem}

This is the central claim of the theory in the strongest form we can prove: the field equation
is gauge covariant for every operator with the stated property. It is a true theorem. It is
also a theorem about a variable. The draft supplies no value for the variable, and the claims
that depend on its particular value, including that it annihilates the Weyl curvature, that
the equation reproduces Einstein's, and everything in Sections~10--12, cannot be stated.

\section{Filling the gap from the literature}\label{sec:fill}

Rule~3 says to repair what can be repaired, so we tried to supply the operator ourselves.

\paragraph{Equivariance is free.} The recipe of (8.1) conjugates a fixed form by the group
element: $\Sh_\eps\xi=[\eps^{-1}\Phi\eps\wedge *\xi]$.

\begin{proposition}[\lean{conjugated\_pairing\_equivariant}]
For any fixed $\Phi$ and any equivariant pairing $B$, the operator
$\Sh_\eps\xi=B(\eps^{-1}\Phi\eps,\,\xi)$ satisfies
$\Sh_{\eps h}(h^{-1}\xi h)=h^{-1}(\Sh_\eps\xi)h$.
\end{proposition}

So the hypothesis of Theorem~\ref{thm:upsilon} holds for every operator of this shape, and the gauge covariance
of the field equation is established for all of them. This is good news for the draft. It also
means that gauge covariance, the property the construction is advertised for, does nothing to
select $\Phi$.

\paragraph{The textbook forms.} The literature has a canonical candidate for the $\Phi_i$: the
tautological Clifford-valued forms
\[
\Phi_1=\sum_a e^a\otimes\gamma_a,\qquad \Phi_2=\sum_{a<b} e^a\wedge e^b\otimes\gamma_a\gamma_b ,
\]
where $\Phi_1$ is the form whose Clifford action defines the Dirac operator
\cite[Chapter~II, Section~5]{LM}. We put these into
the sample operator (9.3) at $\eps=1$ and applied it to the curvature of the spin connection of
a metric, $\xi=\tfrac14 R_{efcd}\,e^e\wedge e^f\otimes\gamma^c\gamma^d$
\cite[Chapter~II, Section~4]{LM}. After taking Hodge
duals, (9.3) has a first term $P_f=\sum_a[\gamma_a,\xi_{af}]$, labelled ``Ricci like'', and a
second term $Q_f=[\gamma_f,S]$ with $S=\sum_{a<b}[\gamma_a\gamma_b,\xi_{ab}]$, labelled
``Ricci scalar like''. The draft's notation allows each bracket to be a commutator or the
symmetric product $\{a,b\}=i(ab+ba)$ of its appendix (the text after (8.1), and (12.27)), and
(9.3) prints plain brackets. We
evaluated all eight combinations on random tensors with the symmetries of a Riemann tensor, in
dimensions $4$ and $6$ with a positive-definite metric
(\lean{experiments/einstein\_check.py}).

\begin{enumerate}
\item In every combination the Weyl part of the curvature is annihilated, as the draft
asserts.
\item With commutators throughout, the second term is identically zero and the operator is the
Ricci tensor alone.
\item In exactly one combination, a commutator in the first term and the symmetric product in
both brackets of the second, the operator is $\mathrm{Ric}-\tfrac12 R\,g$, the Einstein tensor.
The magnitude $\tfrac12$ is the coefficient printed in (9.3). Its sign depends on conventions
for the Hodge star and the Clifford relation that the draft does not state.
\end{enumerate}

The reason for item~2 is short. The Riemann tensor is unchanged when its first pair of indices
is exchanged with its second \cite[Section~1.G]{Besse}, which makes the two orderings in the inner bracket equal, so
their difference vanishes. And had the inner bracket produced the scalar curvature, the outer
bracket would be the commutator of $\gamma_f$ with a scalar, which is also zero. The symmetric
product adds where the commutator subtracts, and
$R_{abcd}\gamma^a\gamma^b\gamma^c\gamma^d=-2R$ supplies the scalar. This identity, and the
identity that makes the first term the Ricci tensor, are the ones used to prove
Lichnerowicz's formula for the square of the Dirac operator \cite[Chapter~II, Section~8]{LM};
their signs depend on the convention for the Clifford relation.

Without its scalar term the equation is $R_{\mu\nu}=-\kappa T_{\mu\nu}$, the generally
covariant equation of Einstein's addendum of 11~November 1915 \cite{Ein15b}, which followed the
paper of 4~November \cite{Ein15a}. It required the hypothesis that the trace of the energy
tensor of matter vanishes, and two weeks later Einstein replaced it with the equation that has
the trace term \cite{Ein15c}; the history is told in \cite{JR}. The reason is the one a modern
reader would give. The divergence of the Ricci tensor is half the gradient of the scalar
curvature, not zero, so the earlier equation is compatible with conservation of energy and
momentum only if the trace of the energy tensor is constant.

\paragraph{What the reconstruction shows.} Read as generously as possible, with forms from the
textbooks and the one bracket assignment that works, the draft's sample operator is correct on
the gravitational part of the curvature. We are glad to be able to report that. But what it is
correct \emph{as} is the Einstein tensor \cite{Ein15c} in Clifford notation, which has been
known since 1915. The curvature in the draft takes
values in all of $u(64,64)$, of which $so(7,7)=\Lambda^2$ is $91$ dimensions out of $16384$.
For the rest, the pair symmetry and Bianchi identity used above are not available, the
literature offers no operator, and the draft does not say what $\Sh$ does. That remainder is
where anything new in the theory would have to be.

\section{Counting the operators}\label{sec:count}

The reconstruction above tried one operator. The draft's description permits many, and the
question of how many has a definite answer.

By the proposition of Section~\ref{sec:fill}, conjugation by $\eps$ makes any fixed pointwise
linear map gauge covariant. What the map must satisfy in order to be defined on the bundle at
all is invariance under the structure group. So the operators of the kind the draft describes
are, at $\eps=1$, the $Spin(7,7)$-equivariant linear maps
\[
\Lambda^2(\mathbb R^{7,7})\otimes u(64,64)\;\longrightarrow\;\Lambda^{13}(\mathbb
R^{7,7})\otimes u(64,64).
\]
These form a real vector space. Its dimension can be computed, because by (8.2)--(8.6)
$u(64,64)$ is the sum of the exterior powers $\Lambda^0,\dots,\Lambda^{14}$ as a representation
(the Clifford algebra and the exterior algebra are isomorphic as representations of the
orthogonal group \cite[Chapter~I, Section~1]{LM}), each exterior power is irreducible except
the middle one, which is a sum of two \cite[Section~19.2]{FH}, and tensor products of
irreducible representations decompose by the formula of Brauer and Klimyk
\cite{Brauer,Klimyk}.

\begin{proposition}\label{prop:count}
\begin{enumerate}
\item The space of such operators has dimension $200$.
\item For each $1\le i\le 13$ the space $[\Lambda^i\otimes u(64,64)]^{Spin(7,7)}$ of
Definition~8.1 has dimension exactly $2$.
\item Every one of the $200$ operators annihilates the Weyl curvature.
\item On the curvature of a metric, the output of any of them is a combination of the
traceless Ricci tensor and the scalar curvature, appearing in Clifford degrees $1$ and $13$.
That is four parameters.
\end{enumerate}
\end{proposition}

Items 3 and 4 follow from Schur's lemma \cite[Section~1.2]{FH} once the decomposition is known.
The curvature tensor of a metric is the sum of three irreducible pieces, the Weyl tensor, the
traceless Ricci tensor and the scalar curvature \cite[Section~1.G]{Besse}. The Weyl piece does
not occur in $\Lambda^{13}\otimes u(64,64)$, and the traceless Ricci tensor and the scalar each
occur there twice. The computation is
\lean{experiments/count\_operators.py}. It is carried out in two independent orders that
agree, the dimensions of the pieces (from Weyl's dimension formula \cite[Section~24.1]{FH}) add
up to the dimensions of the products, and in dimension
$4$ the analogous count ($40$) agrees with a brute-force construction of all the operators as
matrices (\lean{experiments/brute\_force\_d4.py}). It is a computation and not a Lean proof.

Three consequences. First, the draft's statement in its Section~12.4 that the operator ``always
kills off the Weyl curvature'' is true, and it is true of every candidate, so like gauge covariance it selects
nothing. Second, Einstein's equation fixes the ratio of Ricci to scalar, which is a condition
on four of the $200$ parameters; that is the whole of what gravity can say about the operator.
Third, the remaining $196$ parameters govern what the operator does to the part of the
curvature that is not gravitational, and nothing in the draft constrains them.

\section{The action does not give the equation}\label{sec:action}

The draft does say how the operator was chosen: ``for its properties relative to the Bianchi
identity'', within a Lagrangian framework. It gives the action (9.4),
\[
I=\bigl\langle T,\;*\,\Sh_\eps\bigl(F_B+\tfrac12 d_BT+\tfrac13[T,T]\bigr)\bigr\rangle
+\tfrac{\kappa_1}{2}\langle T,*T\rangle ,
\]
modelled on Chern--Simons theory \cite{CS,Witten} (the draft says its model ``abstracts the
Einstein-Hilbert and Chern-Simons actions''), and states in (9.7) that its derivative in the direction of
the connection is the field equation: $\partial_s I(\eps,\vp+s\alpha)=\langle\alpha,\,
\Sh_\eps F_A+\kappa_1{*T}\rangle$. This is a condition on $\Sh$ that can be tested for all
$200$ operators at once.

\begin{proposition}\label{prop:action}
Let $\Sh$ be a $Spin(7,7)$-equivariant operator as in Section~\ref{sec:count}, made gauge
covariant by conjugation. Suppose there is any functional $I(\eps,\vp)$ with
$\partial_sI(\eps,\vp+s\alpha)=\langle\alpha,\Sh_\eps F_A+\kappa_1{*T}\rangle$ for all
$\vp$ and $\alpha$. Then $\Sh=0$.
\end{proposition}

\begin{proof}[Proof sketch]
Write $A=B+T$, so that $F_A=F_B+d_BT+\tfrac12[T\wedge T]$, and encode the operator by its
pairing $K^{abc}(X,Y)=\langle e^c\otimes Y,\,*\Sh(e^a\wedge e^b\otimes X)\rangle$, which is
antisymmetric in $a,b$. An expression is the derivative of a functional only if its own
derivative is symmetric, because second derivatives commute; this is the Helmholtz condition of
the inverse problem of the calculus of variations \cite[Section~5.4]{Olver}. Here that means
$\int K(d_B\beta+[T\wedge\beta],\alpha)=\int K(d_B\alpha+[T\wedge\alpha],\beta)$ for all
$T,\alpha,\beta$. The term without $T$, after an integration by parts, gives
$K^{acb}(Y,X)=-K^{abc}(X,Y)$. The term linear in $T$ has no derivatives and gives
$K^{abc}([Z,X],Y)=K^{acb}([Z,Y],X)$. From the first condition $K^{abc}$ is symmetric in $X,Y$
and totally antisymmetric in $a,b,c$; with the second, each $K^{abc}$ is a symmetric bilinear
form on $u(64,64)$ invariant under the bracket with every $Z$ (\lean{GU.Variational}). The
invariant symmetric forms on $u(64,64)$ are spanned by $\operatorname{Tr}(XY)$ and
$\operatorname{Tr}X\operatorname{Tr}Y$: the algebra is its centre plus $su(64,64)$, whose
complexification $sl(128,\mathbb C)$ is simple, and on a simple Lie algebra the invariant
symmetric forms are the multiples of the Killing form by Schur's lemma
\cite[Section~6]{Humphreys}. So $K^{abc}=v_1^{abc}\operatorname{Tr}(XY)
+v_2^{abc}\operatorname{Tr}X\operatorname{Tr}Y$ with $v_1,v_2$ three-forms. Equivariance
requires $v_1$ and $v_2$ to be fixed by $Spin(7,7)$, and no nonzero three-form is, because
$\Lambda^3$ is an irreducible representation of dimension $364$ \cite[Section~19.2]{FH}.
\end{proof}

The counts agree with the argument. Of the $200$ operators, $48$ satisfy the first condition
and none satisfies both (\lean{experiments/count\_variational.py}). In dimension $4$ the
brute-force construction gives $40$, $8$ and $0$.

The reason is visible in the model. Chern--Simons theory works in three dimensions because
three-dimensional space has an invariant three-form, the volume form, and
$K^{abc}=\epsilon^{abc}\operatorname{Tr}(XY)$ is the solution, which is the Chern--Simons
action \cite{CS,Witten}. Fourteen-dimensional space has
none. The argument uses nothing else about the structure group, so it is unaffected if the
group is reduced to one preserving the splitting of the tangent space of $Y$ into its
four- and ten-dimensional parts, which also fixes no three-form. The second condition also
forces the operator to be invariant under the whole gauge group, so that it cannot depend on
$\eps$. Dependence on $\eps$ is what the ``ship in a bottle'' construction of Section~8 of the
draft consists of.

So (9.4) and (9.8) cannot both stand, whichever operator is chosen. The action (9.4) does have
Euler--Lagrange equations for any $\Sh$, but they contain further terms, the adjoints of
$\Sh\circ d_B$ and of $\Sh\circ[T\wedge\cdot\,]$, that (9.7) omits. And the equation (9.8), for
any nonzero $\Sh$, is not the equation of any action in the connection. The index manipulation
in the proof is checked in Lean. The two facts about invariant forms are standard and are not.
The derivation of the two conditions from (9.7) is ours, done by hand, and has not been
reviewed by anyone.

\section{What would complete it}

The list is now specific. It would suffice to state:
\begin{enumerate}
\item which of the $200$ operators of Section~\ref{sec:count} is meant, or, if it is not among
them, in what way the intended operator differs from the draft's description;
\item a value or a role for $\kappa_1$;
\item the principle that selects it. ``Annihilates the Weyl curvature'' and ``gauge covariant''
hold for all $200$ and cannot be that principle;
\item either an action whose Euler--Lagrange equation is (9.8), or the corrected equation that
the action (9.4) gives, since by Section~\ref{sec:action} these are different.
\end{enumerate}
Item 2 is a line. Items 1, 3 and 4 are the theory.

\paragraph{An offer.} This is meant literally. The signs are fixed, the decomposition is
checked, the gravitational part of the operator is reconstructed, and the candidates are
counted. In the Lean project the
operator is a single field, \lean{FieldEquationInputs.shiab}, waiting for a definition. If the
author can recall the operator, or choose one, it goes in that field, and the claims that
depend on it become statements a proof assistant can check. We would be glad to finish the
job.

\section{Limits of this note}\label{sec:limits}

The results for Sections~5--7 are proved for the abstract data described in Section~1, which
is all the draft's proofs use. The identification of that data with gauge transformations,
one-forms and $h^{-1}\dA h$ is ours and is not machine-checked. The Clifford degree rule used
for (3.27), the isomorphism $Cl(7,7)\cong\mathbb R(128)$, and the two facts about the volume
element used in Section~5 are standard and are taken as known rather than derived. The
bundle-level constructions of Sections~3 and~4 were not formalized; only their pointwise
linear algebra was. The results of Section~\ref{sec:fill} on the Einstein tensor are numerical, for a
positive-definite metric in dimensions $4$ and $6$ rather than signature $(7,7)$ in dimension
$14$, and are not formalized. The counts of Sections~\ref{sec:count} and~\ref{sec:action} are
computed by script for the complexified Lie algebra, which gives the real dimensions; they
count operators that act pointwise and are invariant under the connected group, and requiring
invariance under reflections as well would lower them. Proposition~\ref{prop:action} is the
newest and least tested claim in this note. Its index lemma is machine-checked, its two
representation-theoretic inputs are standard, and the step from the draft's (9.7) to the two
conditions is a hand derivation that no one else has checked. It concerns the variation in
the connection only, as (9.7) does. Nothing here bears on whether a completed theory would
describe nature.

The Lean project accompanying this note builds against Mathlib \texttt{v4.35.0} with Lean~4.35
and contains no \texttt{sorry}. The draft \cite{GU} and the response \cite{NP} were read
directly, and every quotation and number taken from them was checked against the files cited.
The bibliographic details of the papers cited were checked against published records. Books are
cited by chapter or section, for standard material, and those pointers were not checked against
the printed editions.


\section{Conclusion, in plain terms}

Here is what happened, without the notation.

The draft is a recipe. We followed it as far as it goes. The early steps are standard and
correct. In the middle, two steps contradict each other because a minus sign changes partway
through a calculation; we found where, and changed three signs so that everything agrees.
Every long manuscript has errors of this kind and they count for little.

Then the recipe reaches its main step, the one ingredient that is new, and says that the
author used to know what it was and has lost his notes. We tried the nearest thing from the
standard cookbook. Used exactly as the recipe is written, it gives the version of Einstein's
equation that Einstein himself discarded in 1915. Used in the one way that works, it gives
Einstein's equation. Either way it gives nothing new, and it covers a small corner of what
the recipe needs.

Then we counted every ingredient that fits the recipe's description. There are $200$ of them.
Gravity can tell four apart, and those four were settled in 1915. Nothing in the recipe chooses
among the rest. Last, the recipe says where its main equation comes from, and that can be
checked for all $200$ at once. It fails for every one of them. Whichever ingredient is used,
the equation the recipe prints is not the one its own instructions produce.

Four things seem worth drawing from this.

\emph{Wrong and undefined are different conditions.} The sign errors made definite claims, so
they could be tested, found false and corrected. The operator makes no definite
claim, so there is nothing to test. Nguyen and Polya open their conclusions with a remark they attribute to Freeman Dyson: ``It is
better to be wrong than to be vague'' \cite[Section~4]{NP}. A proof assistant enforces that preference mechanically.
It will take a false statement and show that it is false. It cannot do anything with a
statement that has not been made.

\emph{Formalization is a test of completeness, and passing it takes no cleverness.} Nothing in
this note was difficult. The tool's contribution was to decline to continue past a missing
definition, where a sympathetic reader continues without noticing. This is also why
formalizing a theory cannot be delegated to a future, more capable tool. A better prover
checks faster. It does not supply definitions on the author's behalf, and if it did, the
result would be the prover's theory.

\emph{Charity is a method and not only a courtesy.} We granted every standard fact, resolved
every ambiguity in the draft's favor and repaired what could be repaired. What survives that
treatment cannot be attributed to an unsympathetic reading. A gap that remains after the
reader has done everything possible to close it belongs to the text.

\emph{Where the theory can be completed, it returns what is already known.} This is the
observation we would most like to be shown wrong about, and it is easy to show wrong: state
the operator. Until then the claim that the equation unifies anything rests on the part of the
equation that does not yet exist. This is not evidence that the idea is false. It is the
absence of anything that could be evidence either way, and the burden of supplying it rests
with the person making the claim.

A last point concerns this note itself. One of its authors is a language model, which wrote the proofs, and a language model's
assurance that a proof is correct is worth nothing. That is the reason for the form the work
takes. The Lean results do not ask to be believed; anyone can build the project and watch the
kernel accept them. The numerical experiment is a short script and can be rerun. The same standard
that this note applies to the draft is the one it should be held to.

\begin{thebibliography}{99}\small
\bibitem{GU} E.~Weinstein, \emph{Geometric Unity: Author's Working Draft, v~1.0}, April~1,
2021. \url{https://geometricunity.nyc3.digitaloceanspaces.com/Geometric_Unity-Draft-April-1st-2021.pdf}
(retrieved 10~November 2026).
\bibitem{NP} T.~Nguyen and T.~Polya, \emph{A Response to Geometric Unity}, February~23, 2021.
\url{https://files.timothynguyen.org/geometric_unity.pdf} (retrieved 10~November 2026).
\bibitem{Video} E.~Weinstein, \emph{A Portal Special Presentation --- Geometric Unity: A First
Look}, video, uploaded April~2, 2020, \url{https://youtu.be/Z7rd04KzLcg}. Cited here through
the transcription and timestamps in \cite{NP}; we did not consult the video.
\bibitem{Lean4} L.~de~Moura and S.~Ullrich, The Lean~4 theorem prover and programming language,
in \emph{Automated Deduction -- CADE~28}, Lecture Notes in Computer Science 12699, Springer,
2021, pp.~625--635.
\bibitem{Mathlib} The mathlib Community, The Lean mathematical library, in \emph{Proceedings of
the 9th ACM SIGPLAN International Conference on Certified Programs and Proofs (CPP~2020)}, ACM,
2020, pp.~367--381.
\bibitem{KN} S.~Kobayashi and K.~Nomizu, \emph{Foundations of Differential Geometry}, Vol.~I,
Interscience, 1963.
\bibitem{DK} S.~K.~Donaldson and P.~B.~Kronheimer, \emph{The Geometry of Four-Manifolds},
Oxford University Press, 1990.
\bibitem{LM} H.~B.~Lawson, Jr.\ and M.-L.~Michelsohn, \emph{Spin Geometry}, Princeton
University Press, 1989.
\bibitem{Brown} K.~S.~Brown, \emph{Cohomology of Groups}, Graduate Texts in Mathematics 87,
Springer, 1982.
\bibitem{Besse} A.~L.~Besse, \emph{Einstein Manifolds}, Springer, 1987.
\bibitem{FH} W.~Fulton and J.~Harris, \emph{Representation Theory: A First Course}, Graduate
Texts in Mathematics 129, Springer, 1991.
\bibitem{Humphreys} J.~E.~Humphreys, \emph{Introduction to Lie Algebras and Representation
Theory}, Graduate Texts in Mathematics 9, Springer, 1972.
\bibitem{Brauer} R.~Brauer, Sur la multiplication des caract\'eristiques des groupes continus
et semi-simples, \emph{C.~R.~Acad.\ Sci.\ Paris} \textbf{204} (1937), 1784--1786.
\bibitem{Klimyk} A.~U.~Klimyk, Decomposition of the direct product of irreducible
representations of semisimple Lie algebras into irreducible representations (in Russian),
\emph{Ukrain.\ Mat.\ Zh.} \textbf{18} (1966), no.~5, 19--27.
\bibitem{Olver} P.~J.~Olver, \emph{Applications of Lie Groups to Differential Equations},
2nd~ed., Graduate Texts in Mathematics 107, Springer, 1993.
\bibitem{CS} S.-S.~Chern and J.~Simons, Characteristic forms and geometric invariants,
\emph{Ann.\ of Math.} \textbf{99} (1974), 48--69.
\bibitem{Witten} E.~Witten, Quantum field theory and the Jones polynomial, \emph{Comm.\ Math.\
Phys.} \textbf{121} (1989), 351--399.
\bibitem{Ein15a} A.~Einstein, Zur allgemeinen Relativit\"atstheorie, \emph{Sitzungsber.\
K\"onigl.\ Preuss.\ Akad.\ Wiss.\ (Berlin)} (1915), 778--786.
\bibitem{Ein15b} A.~Einstein, Zur allgemeinen Relativit\"atstheorie (Nachtrag), \emph{ibid.}
(1915), 799--801.
\bibitem{Ein15c} A.~Einstein, Die Feldgleichungen der Gravitation, \emph{ibid.} (1915),
844--847.
\bibitem{JR} M.~Janssen and J.~Renn, Arch and scaffold: How Einstein found his field
equations, \emph{Physics Today} \textbf{68} (2015), no.~11, 30.
\end{thebibliography}

\end{document}
